\chapter{Learning the calibration that an entropy certificate needs}
\label{ch:dc-entry}

An inference method is usable only when its calibration assumptions can be
obtained or justified. The preceding detector chapter established when a
complete observation law retains the ideal diffusion information. A finite
experiment must also account for uncertainty in the detector and in the
diffusion scale. This chapter introduces a construction that does so, together
with a precise obstruction to removing all structural assumptions.

\section{The question answered here}

Suppose a bounded signal is recorded through a shutter and a noisy detector.
Can the same finite records both calibrate its local diffusion scale and
certify positive entropy production, without treating either a fitted scale or
a detector-noise estimate as exact?

The answer developed below is conditional on a specified local geometry:
\[
 |f|\leq1,\qquad \Gamma_f=\gamma(1-f^2),\qquad \gamma>0.
\]
The unknown constant $\gamma$ is estimated. The pointwise functional relation
is a premise. Hidden dimension, hidden drift and stationary density are not
fixed to the circle used in the numerical example. Uniform bounds on the
relevant first and second generator derivatives remain additional premises.
This distinction is substantive: a compact countermodel proves that finite
passive records cannot uniformly certify those global bounds over an
unrestricted smooth diffusion class.

\section{What each part of the experiment supplies}

Repeated detector readings of one exposure estimate powers of its underlying
signal through products of distinct readings. An ordinary power of one noisy
reading does not do this. Conditional independence of the detector
randomizations, a common conditional mean and a bounded output range are the
required noise assumptions; the method does not insert a known read-noise
variance.

A constant reference signal and a timed step reference cover a common
multiplicative gain and the width of a uniform shutter. These are actual
additional reference operations, with their own observations and costs.
They do not establish the local-geometry identity or the generator bounds.
A detector with an unknown offset, changing gain or dependent replica noise
requires a different contract.

Time-separated target records provide a finite-step estimate of the average
local diffusivity. The structural relation then converts this into a covered
bound on $\gamma$. A simultaneous region for fixed measured moments carries
its uncertainty through the entropy calculation. Calibration and testing can
use the same target records because their joint coverage is proved; neither
step is silently treated as certain.

\section{A hierarchy with a finite cost}

The polynomial weights
\[
 w_m(y)=1+y^2+\cdots+y^{2m}
\]
obey $w_m(f)\Gamma_f\leq\gamma$. They replace an inverse local-diffusivity
weight by a finite polynomial that independent replicas can measure. Order
$m$ uses $2m+2$ center readings and one reading at each endpoint. All readings,
including additional replicas, consume physical access.

For the stationary cosine-circle example, the ideal limit at order $m$
recovers the fraction
\[
 1-\frac{\binom{2m+2}{m+1}}{4^{m+1}}
\]
of the exact entropy minimum. This fraction tends to one. The constructive
consistency theorem accounts for an unknown diffusion scale, gain calibration,
finite time steps, shrinking exposure and increasing numerical precision.
It does not claim consistency at a fixed nonzero shutter width or practical
optimality of its conservative acquisition schedule.

The sharp result extends beyond a constant hidden current. Let a reversible
hidden process $Z$ modulate angular drift through $c(Z)$, while the measured
signal remains $\cos\phi$. Within the stated ordinary time-reversal class,
the compatible entropy minimum is exactly $\mathbb E[c(Z)^2]/\gamma$.
It can be positive even when $\mathbb E[c(Z)]=0$, and the complete measured
signal law is time-reversible. Thus mean circulation and observed temporal
asymmetry can both vanish while the observation and its local geometry still
force dissipation. This is an analytic extension; the simulation below uses
constant angular drift.

At a finite budget, more replicas and a higher polynomial order can reduce
performance. The implementation therefore retains the lower orders and
compares them under one simultaneous confidence event. Fixed linear contrasts
preserve correlations between measured moments before concentration; remaining
gain uncertainty is paid through the same covered region. Selecting the
largest covered certificate is justified by that event, rather than by an
uncorrected search through individually valid tests.

\section{How to read the extension}

The calibration-limit chapter explains why additional structure is necessary.
The constructive chapter gives the detector contract, the exposure error
bounds, the joint confidence theorem and the hierarchy. The virtual experiment
then evaluates the actual finite procedure, including its simpler comparator
and unsuccessful regimes. An independent audit checks the proofs and the
implementation; the literature chapter identifies the established inference
methods on which the construction depends.

The numerical experiments use a circle to generate records, but the reported
geometry certificates use the broader stated structural class. A fit that
assumes the complete circle dynamics would answer a different inverse problem.
The mathematical results, simulated observations and requirements for a future
physical experiment retain separate evidence status throughout.
